\documentclass[10pt,a4paper]{amsart}
\usepackage{amsmath,amssymb,amsthm,graphicx}
% Retained preamble for the Stacks Project chapter {CH} (shared project preamble).
% The chapter file pinned-original-source.tex only carries \input{{preamble}} on its first line,
% so no preamble text lives inside the chapter itself.  The Stacks Project uses one single shared
% preamble for every chapter: the snapshot stages it three times (sources/stacks-chow/preamble.tex,
% sources/stacks-etale-cohomology/preamble.tex and focused-seeds/AGP/sources/stacks-snapshot-20261011/
% preamble.tex) and those copies are identical in size.  The block below is copied mechanically,
% line for line, from sources/stacks-chow/preamble.tex lines 14-18 and 154-217: the Xy-pic setup and
% the native theorem declarations together with the mathematical macros.  Nothing was added, removed
% or redefined; the wrapper class and the amsmath/amssymb/amsthm/graphicx base packages are supplied
% by the assembler.
% For commutative diagrams we use Xy-pic
% Chapter: stacks-cohomology.tex -- Cohomology of Sheaves.
\usepackage[all]{xy}

% We use 2cell for 2-commutative diagrams.
\xyoption{2cell}

% Theorem environments.
%
\theoremstyle{plain}
\newtheorem{theorem}[subsection]{Theorem}
\newtheorem{proposition}[subsection]{Proposition}
\newtheorem{lemma}[subsection]{Lemma}

\theoremstyle{definition}
\newtheorem{definition}[subsection]{Definition}
\newtheorem{example}[subsection]{Example}
\newtheorem{exercise}[subsection]{Exercise}
\newtheorem{situation}[subsection]{Situation}

\theoremstyle{remark}
\newtheorem{remark}[subsection]{Remark}
\newtheorem{remarks}[subsection]{Remarks}

\numberwithin{equation}{subsection}

% Macros
%
\def\lim{\mathop{\mathrm{lim}}\nolimits}
\def\colim{\mathop{\mathrm{colim}}\nolimits}
\def\Spec{\mathop{\mathrm{Spec}}}
\def\Hom{\mathop{\mathrm{Hom}}\nolimits}
\def\Ext{\mathop{\mathrm{Ext}}\nolimits}
\def\SheafHom{\mathop{\mathcal{H}\!\mathit{om}}\nolimits}
\def\SheafExt{\mathop{\mathcal{E}\!\mathit{xt}}\nolimits}
\def\Sch{\mathit{Sch}}
\def\Mor{\mathop{\mathrm{Mor}}\nolimits}
\def\Ob{\mathop{\mathrm{Ob}}\nolimits}
\def\Sh{\mathop{\mathit{Sh}}\nolimits}
\def\NL{\mathop{N\!L}\nolimits}
\def\CH{\mathop{\mathrm{CH}}\nolimits}
\def\proetale{{pro\text{-}\acute{e}tale}}
\def\etale{{\acute{e}tale}}
\def\QCoh{\mathit{QCoh}}
\def\Ker{\mathop{\mathrm{Ker}}}
\def\Im{\mathop{\mathrm{Im}}}
\def\Coker{\mathop{\mathrm{Coker}}}
\def\Coim{\mathop{\mathrm{Coim}}}

% Boxtimes
%
\DeclareMathSymbol{\boxtimes}{\mathbin}{AMSa}{"02}

%
% Macros for moduli stacks/spaces
%
\def\QCohstack{\mathcal{QC}\!\mathit{oh}}
\def\Cohstack{\mathcal{C}\!\mathit{oh}}
\def\Spacesstack{\mathcal{S}\!\mathit{paces}}
\def\Quotfunctor{\mathrm{Quot}}
\def\Hilbfunctor{\mathrm{Hilb}}
\def\Curvesstack{\mathcal{C}\!\mathit{urves}}
\def\Polarizedstack{\mathcal{P}\!\mathit{olarized}}
\def\Complexesstack{\mathcal{C}\!\mathit{omplexes}}
% \Pic is the operator that assigns to X its picard group, usage \Pic(X)
% \Picardstack_{X/B} denotes the Picard stack of X over B
% \Picardfunctor_{X/B} denotes the Picard functor of X over B
\def\Pic{\mathop{\mathrm{Pic}}\nolimits}
\def\Picardstack{\mathcal{P}\!\mathit{ic}}
\def\Picardfunctor{\mathrm{Pic}}
\def\Deformationcategory{\mathcal{D}\!\mathit{ef}}
\usepackage{cleveref}
\crefname{originallemma}{original lemma}{original lemmas}
\Crefname{originallemma}{Original lemma}{Original lemmas}
\makeatletter
\newlabel{lemma-restrict-direct-image-open}{{08FE}{source}{}{sourcepointer-1}{}}
\newlabel{lemma-restrict-direct-image-open@cref}{{[originallemma][0][]08FE}{[1][0][]source}{}{}{}}
\makeatother
\title{Comparison of an object supported on a closed subset inside an open with its restriction and its extension by zero}
\author{The Stacks Project Authors}
\date{stacks-project:cohomology.tex:local-snapshot-20261011; GNU FDL 1.2 or later}
\begin{document}
\maketitle
\noindent\textit{Editorial note.} All delivered mathematical text is the author's own, copied from the pinned Stacks Project chapter file cohomology.tex; no statement or proof step was rewritten, reordered or removed.\par\medskip
\noindent\textit{Editorial note.} The proof is complete as written and checks both parts by restricting to the two opens; no step is referred to the reader.\par\medskip
\noindent\textit{Editorial note.} The counted spans carry 1 navigation stub and no bibliography citation.\par\medskip
\noindent\textit{Editorial note.} The preamble is the chapter's retained minimal preamble (sources/stacks-cohomology/retained-preamble.tex lines 1-81), a verbatim excerpt of the native preamble.\par\medskip
\noindent\textit{Editorial note.} The label lemma-pushforward-restriction is defined only in cohomology.tex, so no cross-chapter restatement had to be weighed for this card.\par\medskip
\noindent\textit{Allowed standard prerequisites (editorial).}\begin{itemize}
\item The derived pushforward of a restricted object along an open immersion restricts to the same object after intersecting with another open. Same chapter, Lemma lemma-restrict-direct-image-open, invoked at proof lines 8313-8315 and 8317-8320. Supplies the two restriction computations used on V.
\item Extension by zero restricts to extension by zero along an open immersion. Standard, used at proof lines 8321-8322. Supplies the vanishing of the extension by zero on V.
\end{itemize}\medskip
\ifdefined\hypertarget\hypertarget{sourcepointer-1}{}\fi
\noindent\textit{Source reference.} Original lemma 08FE: sources/stacks-cohomology/pinned-original-source.tex, source lines 7901--7901. Pointer outside the counted span, invoked once for context; the load-bearing argument of the counted statement is the author's proof reproduced in full in the proof span.\par\medskip
\begin{lemma}
\label{lemma-pushforward-restriction}
Let $(X, \mathcal{O}_X)$ be a ringed space. Let $j : U \to X$ be an
open subspace. Let $T \subset X$ be a closed subset contained in $U$.
\begin{enumerate}
\item If $E$ is an object of $D(\mathcal{O}_X)$ whose cohomology sheaves
are supported on $T$, then $E \to Rj_*(E|_U)$ is an isomorphism.
\item If $F$ is an object of $D(\mathcal{O}_U)$ whose cohomology sheaves
are supported on $T$, then $j_!F \to Rj_*F$ is an isomorphism.
\end{enumerate}
\end{lemma}

\begin{proof}
Let $V = X \setminus T$ and $W = U \cap V$. Note that $X = U \cup V$ is an
open covering of $X$. Denote $j_W : W \to V$ the open immersion.
Let $E$ be an object of $D(\mathcal{O}_X)$ whose cohomology sheaves are
supported on $T$. By
Lemma \ref{lemma-restrict-direct-image-open} we have
$(Rj_*E|_U)|_V = Rj_{W, *}(E|_W) = 0$ because $E|_W = 0$ by our assumption.
On the other hand, $Rj_*(E|_U)|_U = E|_U$. Thus (1) is clear.
Let $F$ be an object of $D(\mathcal{O}_U)$ whose cohomology sheaves
are supported on $T$. By
Lemma \ref{lemma-restrict-direct-image-open} we have
$(Rj_*F)|_V = Rj_{W, *}(F|_W) = 0$ because $F|_W = 0$ by our assumption.
We also have $(j_!F)|_V = j_{W!}(F|_W) = 0$ (the first equality is immediate
from the definition of extension by zero). Since both
$(Rj_*F)|_U = F$ and $(j_!F)|_U = F$ we see that (2) holds.
\end{proof}

\end{document}
